\section{Detailed proofs of the Galois-group and specialization steps}
\label{app:detailed-proofs}

\subsection{The full joint group}\label{app:joint-group}
We expand the group argument in Proposition~\ref{prop:product}.
All groups are reduced complex linear algebraic groups.

\begin{lemma}\label{lem:app-group-interface}
Let $\phi:G\to H$ be a group morphism.
Its image is closed, $G^0$ is normal of finite index, and
\begin{equation}\label{eq:app-group-interface}
\begin{gathered}
 \phi(G^0)=\phi(G)^0,\qquad \dim G=\dim\Lie G,\\
 \Lie(\ker\phi)=\ker d\phi,\qquad
 \Lie(\phi(G))=d\phi(\Lie G),\\
 \Lie(A\cap B)=\Lie A\cap\Lie B
 \quad(A,B\subset G\text{ closed}).
\end{gathered}
\end{equation}
If $A$ is connected, then $\Lie A\subset\Lie B$ implies $A\subset B$.
A closed subgroup of a connected group with equal dimension is the whole group.
\end{lemma}
\begin{proof}
Closed images and identity components follow from
\cite[Proposition~1.2(b), p.~46; Corollary~1.4(a)--(b), p.~47]{Borel1991}.
Smoothness identifies group dimension with tangent-space dimension
\cite[Proposition~1.2(a), p.~46]{Borel1991}
\cite[\S9.1, p.~65]{Humphreys1975}.
The kernel and intersection identities are
\cite[\S12.5, Theorem, p.~85]{Humphreys1975}.
The image identity uses characteristic-zero separability
\cite[\S7.2, proof, p.~105]{Borel1991}.
If $\Lie A\subset\Lie B$, then $A\cap B$ has dimension $\dim A$.
Connected algebraic groups are irreducible
\cite[Proposition~1.2 and following paragraph, \mbox{pp.}~46--47]{Borel1991}.
Proper closed irreducible subsets have smaller dimension
\cite[\S3.2, Proposition, p.~25]{Humphreys1975}.
Apply this to $(A\cap B)^0\subset A$.
It gives $A\subset B$; see also
\cite[\S7.1(2), p.~105]{Borel1991}
\cite[\S13.1, Theorem, p.~87]{Humphreys1975}.
The last assertion follows by applying the same argument to $B^0\subset G$.
\end{proof}

\begin{lemma}\label{lem:app-pv-compositum}
Let $E/K$ and $F/K$ be PV extensions with algebraically closed constants $\C$.
Assume their compositum $M=EF$ has the same constants.
Then $M/K$ and $M/F$ are PV, and restriction gives
\begin{equation}\label{eq:app-pv-restriction}
 \Gal(M/F)\simeq\Gal(E/(E\cap F)).
\end{equation}
The right side is normal in $\Gal(E/K)$.
Its quotient is a quotient of $\Gal(F/K)$.
The intersection $E\cap F$ is PV over $K$.
The maps $\Gal(M/K)\to\Gal(E/K),\Gal(F/K)$ are onto.
\end{lemma}
\begin{proof}
The block sum of fundamental matrices generates $M/K$.
The matrix for $E$ also generates $M/F$.
Both subfields $E,F$ are invariant under $\Gal(M/K)$.
PV Galois correspondence gives the two surjections
\cite[Proposition~1.34(1)--(2), p.~26]{VanDerPutSinger2002Draft}.
The restriction $\Gal(M/F)\to\Gal(E/K)$ is injective: an automorphism
in its kernel fixes $E$ and $F$, hence fixes $M=EF$, so it is the identity.
Put $H=\operatorname{im}(\Gal(M/F)\to\Gal(E/K))$.
It is closed by Lemma~\ref{lem:app-group-interface}.
Surjectivity and conjugation show that $H$ is normal.
Its fixed field is $E\cap F$, proving
\eqref{eq:app-pv-restriction} by the same correspondence.
Since $H$ is normal, the PV correspondence also makes $E\cap F$ PV over $K$
\cite[Proposition~1.34(2), p.~26]{VanDerPutSinger2002Draft}.
The map to $\Gal(E/K)/H$ kills $\Gal(M/F)$.
It therefore factors through $\Gal(F/K)$.
\end{proof}

\begin{proof}[Detailed proof of Proposition~\ref{prop:product}]
\emph{1. Base change.}
Put $V=\mathsf G_{u,\mathcal U}$ and $H_0=\Gal(QT'/T')$.
The matrices in \eqref{eq:core-unipotent-matrix} are lower unitriangular.
Thus $V$ is a closed unitriangular subgroup, hence nilpotent.
Every nonidentity unipotent has connected cyclic closure
\cite[\S2.1, p.~56; \S7.3, Remark, p.~106]{Borel1991}.
Each element therefore belongs to $V^0$.
Hence $V$ is connected and solvable.

The generators of $T'$ satisfy
\[
 t_h'=\lambda_ht_h,\qquad
 (z^{1/\ell})'=\frac{z^{1/\ell}}{\ell z}.
\]
Thus $T'/k$ is PV with abelian group.
For each block field $\mathcal P_i$, Lemma~\ref{lem:app-pv-compositum}
identifies $\Gal(\mathcal P_iT'/T')$ with a normal subgroup of
$\Gal(\mathcal P_i/k)$ having abelian quotient.
The block group contains $\SL_2$ by Lemma~\ref{lem:groups}.
Since $\SL_2=[\SL_2,\SL_2]$, the normal subgroup contains $\SL_2$.
The group over $T'$ fixes the block Wronskian, which belongs to $T'$:
\[
 W_i\in\C^\times\{z^{-1},\ z^{-2\nu-1},\ e^{sz}/z\}.
\]
Thus the base-changed block group equals $\SL_2$.
All composita have constants $\C$, by the standing constant-field hypothesis.

\emph{2. Projective independence.}
From \eqref{eq:core-block-transport},
\[
 w_i\mathbf1=Y_{\mathfrak c}^{-1}B_i^{-1}Y_iC_i^{-1}
 \quad\Longrightarrow\quad w_i\in Q\cap T'.
\]
Thus \eqref{eq:core-character} holds on the full group $G_Q$.
It gives $\psi(G_Q)\subset\bar P_{\mathrm{SL}}=\psi(P_{\mathrm{SL}})$,
where $P_{\mathrm{SL}}$ is the subgroup denoted $P$ in the body, with its conjugated block embedding.
Each representative projection is onto $\PGL_2$.
Each block field is a PV subextension of $Q/k$.
Lemma~\ref{lem:app-group-interface} makes
$\Lie(\psi(G_Q)^0)$ subdirect in the representative factors.

We use an elementary consequence of Lemma~\ref{lem:goursat}.
A proper subdirect Lie algebra in $\mathfrak{sl}_2^r$ has a graph projection
onto some pair of factors.
For $r=2$, use Lemma~\ref{lem:goursat}.
Proceed by induction on $r$.
If the first $r-1$ projection is proper, apply induction.
Otherwise it is $\mathfrak s=\mathfrak{sl}_2^{r-1}$.
For $\mathfrak h\subset\mathfrak s\oplus\mathfrak{sl}_2$, put
$\mathfrak j=\{b:(0,b)\in\mathfrak h\}$.
It is an ideal in $\mathfrak{sl}_2$.
If $\mathfrak j=\mathfrak{sl}_2$, subdirectness gives the full sum.
Otherwise $\mathfrak j=0$, and $\mathfrak h$ is the graph of a surjection
$\mathfrak s\to\mathfrak{sl}_2$.
The factor images are commuting ideals.
Exactly one is nonzero, giving the required pair.
The standard brackets $[h,e]=2e$, $[h,f]=-2f$, and $[e,f]=h$
show that $\mathfrak{sl}_2$ is simple.
The calculation in (F4) shows that its automorphisms are inner.

A graph pair corresponds to a rank-one twist by Lemma~\ref{lem:twist}.
For completeness, let $G_Q=\Gal(Q/k)$ be the full joint group used in descent.
The projective image of $G_Q^0$ in the $i,j$ factors is the graph group by
Lemma~\ref{lem:app-group-interface}.
Thus $\pi_j=\chi C_0\pi_iC_0^{-1}$ on $G_Q^0$.
On $\operatorname{Hom}(V_i,V_j)$, the line $\C C_0$ is the unique invariant line:
the other summand is the irreducible adjoint representation.
Its invariant subspaces are ideals in $\mathfrak{sl}_2$.
Normality of $G_Q^0$ makes this line $G_Q$-stable.
PV descent gives a rational basis $B_0$ of that line.
In the Hom system it has the form
\[
 B_0=vY_jC_0Y_i^{-1},\qquad
 B_0'+B_0M_i-M_jB_0=rB_0,\qquad r=v'/v\in k.
\]
Here $C_0\in\GL_2(\C)$ is the graph matrix, and $M_i,M_j$ are the block
system matrices.
Thus $B_0$ is invertible.
A scalar solution of $u'/u=-r$ turns it into an intertwiner.
Distinct representatives belong to distinct linkage classes.
They therefore have no graph pair.
The induction gives the full direct-sum Lie algebra.
The dimension criterion in Lemma~\ref{lem:app-group-interface} now yields
\begin{equation}\label{eq:app-projective-full}
 \psi(G_Q^0)=\psi(G_Q)=\bar P_{\mathrm{SL}}.
\end{equation}

\emph{3. The full group after base change.}
Wronskians and \eqref{eq:core-character} give $H_0\subset P_{\mathrm{SL}}$.
Restriction $\Gal(QT'/k)\to G_Q$ is onto.
Lift any $g,h\in G_Q$ to this group.
Their restrictions to $T'$ commute.
Their commutator therefore fixes $T'$, giving $[g,h]\in H_0$.
For $x,y\in P_{\mathrm{SL}}$, equation~\eqref{eq:app-projective-full} supplies $g,h\in G_Q$
with the same respective projective images.
Thus $g=xz_1$, $h=yz_2$, where $z_1,z_2\in\ker\psi$ are central.
Hence $[x,y]=[g,h]\in H_0$.

Put $u_+(b)=\left(\begin{smallmatrix}1&b\\0&1\end{smallmatrix}\right)$,
$u_-(b)=\left(\begin{smallmatrix}1&0\\b&1\end{smallmatrix}\right)$ and
$d(t)=\operatorname{diag}(t,t^{-1})$.
Use the commutator convention $[x,y]=xyx^{-1}y^{-1}$.
For $t^2\ne1$,
\[
 u_+(b)=[d(t),u_+(b/(t^2-1))],\qquad
 u_-(b)=[d(t),u_-(b/(t^{-2}-1))].
\]
Row elimination shows that these elementary matrices generate $\SL_2(\C)$.
Thus $P_{\mathrm{SL}}$ is perfect, and $P_{\mathrm{SL}}\subset H_0$.
Consequently $H_0=P_{\mathrm{SL}}$, including all finite components.

\emph{4. The unipotent factor.}
Both $\mathcal U/T'$ and $QT'/T'$ are PV.
Lemma~\ref{lem:app-pv-compositum} gives surjective restriction maps from
$\Gal(\Omega/T')$ to both groups.
Their product is injective since $\Omega=\mathcal U(QT')$.
Write its image as $H\subset V\times P_{\mathrm{SL}}$ and put
\[
 H_V=\{v:(v,1)\in H\},\qquad H_P=\{p:(1,p)\in H\}.
\]
Conjugation by lifts proves normality in $V,P_{\mathrm{SL}}$, respectively.
Multiplication of lifts gives the isomorphism
\[
 V/H_V\longrightarrow P_{\mathrm{SL}}/H_P,
 \qquad vH_V\longmapsto pH_P
 \quad((v,p)\in H).
\]
Two choices of $p$ differ by $H_P$. The map
$V\to P_{\mathrm{SL}}/H_P$ has kernel $H_V$, so it induces the displayed
isomorphism.
The common quotient is solvable and perfect, hence trivial.
Thus $H=V\times P_{\mathrm{SL}}$, proving \eqref{eq:core-product}.
Also $\mathcal U\cap QT'=T'$: its PV group is the same common quotient,
by Lemma~\ref{lem:app-pv-compositum}.
\end{proof}


\subsection{The tower induction}\label{app:tower-induction}
Keep the hypotheses of Lemma~\ref{lem:tower}.
The nodes $\eta_1,\ldots,\eta_m$ are distinct and nonzero.
The heights $h_j$ are positive integers.
Every rate in $\mathfrak S$ belongs to $\Gamma$.
The finite support contains $(0,0)$ and every $(-\eta_j,0)$.
Use $\Phi_{0,0}=L$ and $\Phi_{-\eta_j,0}=G_{\eta_j,1}$,
as in \eqref{eq:core-G}.
Put
\begin{equation}\label{eq:appendix-tower-assertions}
\begin{gathered}
 H(h)=\sum_j(h_j-1),\qquad
 \mathcal T_h(\mathbf p)=\sum_jp_j(L)G_{\eta_j,h_j}/z,\\
 \trdeg_T\mathcal N_h=\#\mathfrak S+H(h),\qquad
 \mathcal N_h^D=\C,\\
 \mathcal T_h(\mathbf p)\notin D\mathcal N_h
 \quad(0\ne\mathbf p\in\C[L]^m).
\end{gathered}
\end{equation}
The degree assertion means independence of all displayed generators.
We prove these assertions by induction on $H(h)$.

\begin{lemma}\label{lem:appendix-primitive-extension}
Let $K$ be a differential field of characteristic zero.
Suppose $K^D=\C$, $a\in K\setminus DK$ and $Du=a$.
Then $u$ is transcendental over $K$, and $K(u)^D=\C$.
\end{lemma}
\begin{proof}
If $u$ were algebraic, put $n=[K(u):K]$.
Derivation commutes with trace, as in Lemma~\ref{lem:primitive}.
Thus
\[
 D\bigl(n^{-1}\operatorname{Tr}_{K(u)/K}u\bigr)=a,
\]
contrary to $a\notin DK$.
For $Df=0$, write $f=P(u)/Q(u)$ in lowest terms.
Take $Q$ monic.
The identity $QDP=PDQ$ gives $Q\mid DQ$.
Monicity gives $\deg_u DQ<\deg_u Q$, so $DQ=0$.
If $\deg_u Q=r>0$, its next coefficient gives
$Dq_{r-1}+ra=0$.
Hence $Q=1$.
If $\deg_u P=s>0$, then $p_s\in\C^\times$ and
$Dp_{s-1}+sp_sa=0$.
Both equations contradict $a\notin DK$.
Thus $f\in\C$.
\end{proof}

For independent variables $Y_i$ over $K$, put $\partial_i=\partial/\partial Y_i$.
Their common kernel in $K(Y_i)$ is $K$.
A reduced numerator and denominator must divide their partial derivatives.
Degree comparison makes those derivatives zero.
If $DY_i\in K$, then $[D,\partial_i]=0$.
This follows on $K$ and on each $Y_i$.

\emph{1. Height zero.}
The vectors $\operatorname{res}(\Phi_\sigma)$ are nonzero multiples
of distinct unit vectors.
Lemma~\ref{lem:primitive} gives independence and constant field $\C$.
Put $\mathfrak U=\{(-\eta_j,0):1\le j\le m\}$ and
$\mathfrak S_0=\mathfrak S\setminus\{(0,0)\}$.
Suppose $Dy=\mathcal T_{\mathbf1}(\mathbf p)$.
Take partial derivatives over $T(L)$ in the remaining $\Phi_\sigma$.
They commute with $D$.
The constant-field assertion gives
\[
 \partial_{-\eta_j,0}y=Q_j(L),\quad
 \frac{dQ_j}{dL}=p_j,\qquad
 \partial_\sigma y=c_\sigma\in\C
 \quad(\sigma\in\mathfrak S_0\setminus\mathfrak U).
\]
Thus
\[
\begin{aligned}
 y_0&=y-\sum_jQ_j(L)G_{\eta_j,1}
       -\sum_{\sigma\in\mathfrak S_0\setminus\mathfrak U}
           c_\sigma\Phi_\sigma\in T(L),\\
 Dy_0&=-\sum_jQ_j(L)e^{-\eta_jz}/z
       -\sum_{\sigma\in\mathfrak S_0\setminus\mathfrak U}
           c_\sigma\phi_\sigma.
\end{aligned}
\]
Lemma~\ref{lem:reduction} gives $y_0\in R_1$.
Choose $j$ with $p_j\ne0$.
The coefficient of weight $-\eta_j$ satisfies
\[
 Du-\eta_ju=-Q_j(L)/z+g(z),\qquad
 u\in\C(z,L),\quad g\in\C(z).
\]
Here $Q_j$ is nonconstant.
Distinct nodes give distinct weights.
Lemma~\ref{lem:twisted-obstruction} excludes this equation.

\emph{2. Independence and constants at positive height.}
Assume \eqref{eq:appendix-tower-assertions} at every smaller height.
Choose $j$ with $h_j\ge2$.
Let $e_j$ be the $j$th unit vector.
Put $\widetilde h=h-e_j$.
For $K=\mathcal N_{\widetilde h}$, its obstruction gives
\[
 a=G_{\eta_j,h_j-1}/z\notin DK.
\]
Use coefficient $1$ at $j$ and $0$ elsewhere.
Apply Lemma~\ref{lem:appendix-primitive-extension} to $Y=G_{\eta_j,h_j}$.
Since $DY=a$, independence and constants follow at $h$.
The fixed meromorphic domain also gives the constant assertion directly.

\emph{3. The obstruction at positive height.}
Put $\mathcal J_h=\{j:h_j\ge2\}$ and $h^-_j=\max(1,h_j-1)$.
The independence just proved gives
\[
 \mathcal N_h=\mathcal N_{h^-}(Y_j:j\in\mathcal J_h),\qquad
 Y_j=G_{\eta_j,h_j},\quad DY_j\in\mathcal N_{h^-}.
\]
Suppose $Dy=\mathcal T_h(\mathbf p)$ with $\mathbf p\ne0$.
Take partials in the $Y_j$ for $j\in\mathcal J_h$, fixing $\mathcal N_{h^-}$.
Their commutation with $D$ and the constant assertion give
\[
 \partial_jy=Q_j(L),\quad \frac{dQ_j}{dL}=p_j
 \quad(j\in\mathcal J_h),\qquad
 y_0=y-\sum_{j\in\mathcal J_h}Q_j(L)Y_j\in\mathcal N_{h^-}.
\]
Differentiation gives
\begin{equation}\label{eq:appendix-tower-descent}
 Dy_0=\mathcal T_{h^-}(R),\qquad
 R_j=\begin{cases}
 p_j,&h_j=1,\\
 -Q_j,&h_j\ge2.
 \end{cases}
\end{equation}
A nonzero $p_j$ at height one survives.
Otherwise some $Q_j$ has nonzero derivative.
Thus $R\ne0$, while $H(h^-)<H(h)$.
The induction hypothesis excludes \eqref{eq:appendix-tower-descent}.

\emph{4. The listed solvable coordinates.}
Enlarge $\mathfrak S$ to contain their first-level supports.
Put $\mathcal P=T(\Phi_\sigma:\sigma\in\mathfrak S)$.
Choose a listed residue basis $B_1,\ldots,B_\rho$.
Lemma~\ref{lem:primitive} gives, for every other first-level coordinate,
\[
 F_a=\sum_i\kappa_{ai}B_i+r_a+c_a,\qquad
 \kappa_{ai},c_a\in\C,\quad r_a\in R_0.
\]
The higher $G_{\eta,k}$ are independent over $\mathcal P$.
Since $L\in\mathcal P$, so are $\Ein_k(\eta z)=L^k/k!-G_{\eta,k}$.
They remain independent over $T(B_1,\ldots,B_\rho)\subset\mathcal P$.
Eliminating the affine graphs leaves a polynomial ring in the $B_i$
and the $\tau$ distinct higher coordinates.
Hence the degree is $\rho+\tau$ and these graphs generate the ideal.


\subsection{The relation ideal}\label{app:ideal-details}
Retain the hypotheses of Lemma~\ref{lem:ideal}.
Put
\[
 \mathcal R=\C(z)[\Theta^{\pm1}],\qquad n=\rho+\tau,
 \qquad m=\#\mathcal C.
\]
Let $\mathcal I$ be the ideal of the stated generators.
The duplicate, constant, exponential, residue and linkage relations
are monic graphs. Successive substitution eliminates each dependent
variable. The fourth-coordinate shifts in \eqref{eq:core-quadric}
then give the following quotient.
\begin{equation}\label{eq:app-ideal-model}
 B_{\mathcal R}=
 \frac{\mathcal R[Z_1,\ldots,Z_n,(X_{\mathfrak c},Y_{\mathfrak c},
 S_{\mathfrak c},W_{\mathfrak c})_{\mathfrak c\in\mathcal C}]}
 {(X_{\mathfrak c}W_{\mathfrak c}-Y_{\mathfrak c}S_{\mathfrak c}
       +v_{\mathfrak c}:\mathfrak c\in\mathcal C)},
 \qquad v_{\mathfrak c}=u_{\mathfrak c}/z\in\mathcal R^\times.
\end{equation}
Here $W_{\mathfrak c}$ is the shifted fourth variable.
The variable $X_{\mathfrak c}$ represents $y_{\mathfrak c}$.

Lemma~\ref{lem:quadric} gives an $\mathcal R$-basis: use the $Z$-monomials
and, in each block, monomials avoiding $X_{\mathfrak c}W_{\mathfrak c}$.
Thus $B_{\mathcal R}$ is free over $\mathcal R$. It injects into its generic fibre.
That fibre is geometrically integral by Lemma~\ref{lem:quadric}
and \cite[Proposition~(4.6.5)(ii), p.~69]{EGAIV2}.
Hence $B_{\mathcal R}$ is a domain.

Put $M=\prod_{\mathfrak c}X_{\mathfrak c}$.
The same basis shows $M\ne0$. Solving each determinant gives
\begin{equation}\label{eq:app-ideal-chart}
 B_{\mathcal R}[M^{-1}]
 =\mathcal R[Z_1,\ldots,Z_n,
      (X_{\mathfrak c}^{\pm1},Y_{\mathfrak c},S_{\mathfrak c})_
        {\mathfrak c\in\mathcal C}],
 \qquad
 W_{\mathfrak c}=\frac{Y_{\mathfrak c}S_{\mathfrak c}-v_{\mathfrak c}}
 {X_{\mathfrak c}}.
\end{equation}
The earlier degree proof makes these selected functions independent over
$T$. Lemma~\ref{lem:exponential-log} identifies
$\operatorname{Frac}\mathcal R$ with $T$.
Evaluation on \eqref{eq:app-ideal-chart} is therefore injective.
Its restriction to $B_{\mathcal R}$ is injective too.
Consequently $\mathcal I$ is the full Laurent kernel and is prime.
Localization in $\mathcal R\setminus\{0\}$ gives the kernel over $T$.

Let $k_1\subset\C$ contain the generator data and residue constants.
The same substitutions and monic divisions work over
$\mathcal R_1=k_1(z)[\Theta^{\pm1}]$.
Their monomial basis maps to the identical basis over $\mathcal R$.
Thus $B_{\mathcal R_1}\hookrightarrow B_{\mathcal R}$.
Evaluation remains injective, proving descent of the full kernel.
Its fraction field is rational in the selected variables over
$k_1(z)(\Theta)$. Its transcendence degree over $k_1(z)$ is $d+n+3m$.

For ordinary polynomial variables, adjoin $V_h$ with
$\Theta_hV_h-1=0$. Replace negative powers by $V_h$ and eliminate
$\Theta,V$. The quotient is the Laurent quotient above.
Hence contraction gives the full coordinate kernel.
Equivalently, clear negative powers, saturate by
$\prod_h\Theta_h$, and eliminate.
Listed exponential variables require the corresponding saturation too.

\subsection{Specialization at nonzero algebraic points}
\label{app:specialization-details}
Assume the hypotheses of Theorem~\ref{thm:values}.
Thus the list is closed, normalized and saturated.
Its rates, kernels, scales and constants are algebraic.
Its Bessel orders are rational; its Kummer parameters are rational
and nonintegral. Scales are nonzero.
The primitives satisfy $(\mathrm{B2}^E)$, including $F(0)=0$.
First replace integer-order blocks by their order-zero blocks.
Put
\[
 R=\Qbar[z^{\pm1}][\Theta^{\pm1}],\qquad
 \mathcal R_\alpha=\Qbar[\Theta^{\pm1}],\qquad
 n_{\mathrm{val}}=d+\rho^E+\tau+3m,\quad m=\#\mathcal C.
\]

The residue vectors have algebraic entries. Hence their basis
coefficients $\kappa_{ji}$ are algebraic.
For $\mu\in\Qbar$, put $D_\mu g=g'+\mu g$.
On $\Qbar[z^{\pm1}]$,
\begin{equation}\label{eq:app-laurent-image}
 \operatorname{im}D_\mu
 =\ker\bigl(q\mapsto\Res_0 e^{\mu z}q(z)\,dz\bigr).
\end{equation}
The residue vanishes on $D_\mu g$ since
$e^{\mu z}D_\mu g=(e^{\mu z}g)'$.
For $\mu=0$, differentiate Laurent monomials.
For $\mu\ne0$, $D_\mu$ is invertible on polynomials.
Its negative-power reduction is
$z^{-j-1}\equiv(\mu/j)z^{-j}$ for $j\ge1$.
The remaining class is $z^{-1}$, with residue $1$.
This proves \eqref{eq:app-laurent-image}.
Apply it to each residue-zero first-level relation.
It gives $r_j\in R$ in \eqref{eq:residue-relations}.
Zero-origin normalization gives
\[
 c_j=-\lim_{z\to0}r_j(z,e^{\lambda_1z},\ldots,e^{\lambda_dz})
     \in\Qbar.
\]
The limit exists because the corresponding primitive difference is entire.
It is a finite Laurent--exponential constant coefficient.

All other graph coefficients lie in $R$.
For the shift $(f_\nu,f'_\nu)\mapsto(f_{\nu+1},f'_{\nu+1})$ and the shift
$(f_{-\nu},f'_{-\nu})\mapsto(f_{-\nu-1},f'_{-\nu-1})$, respectively,
the solution-coordinate determinants are
\[
 -\frac{4(\nu+1)^2}{c^2z^2},\qquad
 -\frac{c^2z^2}{4\nu^2}.
\]
Their product is $(\nu+1)^2/\nu^2$.
The normalized noninteger shifts have $\nu(\nu+1)\ne0$.
Kummer contiguity has determinant $1$ on each diagonal block.
The exponential links have unit determinants in $R$.
Integer-order reduction has determinant $(-1)^n$ on each diagonal block
by Lemma~\ref{lem:integer}.
Scale and order signs use identities or permutations.
Thus every required map and inverse specializes at $\alpha\ne0$.

Appendix~\ref{app:arithmetic} supplies the $E$-function bounds.
Use the vector in \eqref{eq:value-system} after order-zero reduction.
Its primitive equations have Laurent-polynomial coefficients.
Its tower equations are \eqref{eq:ein-tower}.
The block equations are
\[
\begin{gathered}
 I''=c^2I-I'/z,\qquad
 S''=c^2S-S'/z+2I'/z,\\
 F''=(s-z^{-1})F'+asF/z,\\
 \Sigma''=(s-z^{-1})\Sigma'+as\Sigma/z+(2F'-sF)/z,\\
 f''_{\nu,c}=c^2f_{\nu,c}-(2\nu+1)f'_{\nu,c}/z.
\end{gathered}
\]
Consequently $f$ satisfies a first-order system with coefficient matrix
in $\Qbar[z,z^{-1}]$.
Every $\alpha\in\Qbar^\times$ is ordinary.
Each component of $f$ has algebraic Taylor coefficients. Any polynomial
relation over $\C(z)$ therefore descends to $\Qbar(z)$ by solving its
coefficient equations over $\Qbar$; the two transcendence degrees agree.
Lemma~\ref{lem:ideal} and the generic degree formula give
\begin{equation}\label{eq:app-specialized-degree}
 \trdeg_{\Qbar}\Qbar(f(\alpha))
 =\trdeg_{\Qbar(z)}\Qbar(z)(f)=n_{\mathrm{val}}
\end{equation}
by \cite[Theorem~1.1, p.~370]{Beukers2006}.
The extra system exponentials are Laurent monomials in the basis
exponentials. Thus they leave both fields unchanged.

Specialize the graph eliminations at $z=\alpha$.
They give \eqref{eq:app-ideal-model} over $\mathcal R_\alpha$, with
$n=\rho^E+\tau$ and
\[
 v_{\mathfrak c}\in
 \{\alpha^{-1},\Theta^{v_s}\alpha^{-1},2\nu\alpha^{-1}\}
 \subset \mathcal R_\alpha^\times.
\]
Write $B_\alpha$ for this quotient.
The free-basis argument above makes $B_\alpha$ a domain.
Its generators vanish on the values.

Let $s_0$ count the blocks with $X_{\mathfrak c}(\alpha)=0$.
In each such block,
\[
 Y_{\mathfrak c}(\alpha)S_{\mathfrak c}(\alpha)
   =v_{\mathfrak c}(e^{\lambda_1\alpha},\ldots,e^{\lambda_d\alpha})
   \ne0.
\]
Over the exponential-value field, this block needs two generators:
$Y_{\mathfrak c}(\alpha),W_{\mathfrak c}(\alpha)$.
Every other block needs three, by its determinant.
All values are therefore generated by at most $n_{\mathrm{val}}-s_0$ numbers.
Equation~\eqref{eq:app-specialized-degree} forces $s_0=0$.
It then makes the $n_{\mathrm{val}}$ selected values, including basis exponentials,
algebraically independent.

Evaluation now extends to $B_\alpha[M^{-1}]$.
By \eqref{eq:app-ideal-chart}, this is a Laurent polynomial ring
in those selected variables. Its evaluation is injective.
Since $B_\alpha\hookrightarrow B_\alpha[M^{-1}]$, evaluation on
$B_\alpha$ is injective. Hence $I_\alpha=\mathfrak P_\alpha$.
In particular $y_{\mathfrak c}(\alpha)\ne0$.
Nonvanishing was established before localization.
The system denominators are powers of $z$; all coordinate maps have
unit determinants in $R$. These facts account for every excluded point.

For several points put $\widetilde q_i(z)=q_i(\alpha_i z)$.
Primitive kernels become $\alpha_iq_{i,\mu}(\alpha_i z)$.
Residues retain their values under
$(\mu,p)\mapsto(\alpha_i\mu,p/\alpha_i)$.
Tower pairs become $(\alpha_i\eta_i,k)$.
Derivative coordinates acquire the nonzero factors $\alpha_i$.
Close and normalize the combined list; saturate its rates.
Test residues, tower pairs and linkage classes on this list.
The one-point result at $1$ gives the several-point assertions.

